\section{Hermitian geometry of Kato manifolds}\label{metrici}

Kato manifolds cannot admit Kähler metrics since their fundamental group is isomorphic to $\mathbb{Z}$. However, a large class of them carry \emph{locally conformally Kähler metrics (lcK)}, as shown in \cite[Th.\,2.2]{iop}. In this section, we give a characterization for the existence of lcK metrics on Kato manifolds. Furthermore, we investigate the possibility of endowing a Kato manifold with other special Hermitian metrics, such as balanced, pluriclosed, Hermitian symplectic or strongly Gauduchon.

Recall the following definition:
\begin{definition}\label{deflck} A \emph{locally conformally Kähler metric (lcK)} on a complex manifold~$X$ is a Hermitian metric $\Omega$ satisfying $d\Omega=\theta\wedge\Omega$ for a closed one-form $\theta$. If $\theta$ is not exact, $\Omega$ is called strictly lcK.
\end{definition}

This definition is equivalent to any of the following formulations:
\begin{enumerate}
\item There exists a covering with open sets $(U_i)_i$ of $X$ and smooth functions $f_i$ on every $U_i$ such that $e^{-f_i}\Omega$ is Kähler.
\item The universal cover $\wt{X}$ admits a Kähler metric $\wt{\Omega}$ such that $\gamma^*\wt{\Omega}=c_{\gamma}\wt{\Omega}$, $c_{\gamma}>0$, for any deck-transformation $\gamma$.
\end{enumerate}

For our characterization purpose, we start by giving a general principle used to modify Kähler metrics on modifications of balls, already used in \cite{bru} and \cite{iop}, which will be needed.

\begin{lemma}\label{extKA}
Let $\pi:\hat\CC^n\to \CC^n$ be a proper modification at $0$, let $\rho(z)=\rho(\|z\|)\in\ce(\CC^n,\RR)$ be a strictly plurisubharmonic function depending only on $\|z\|$, and let $\omega$ be a \Ka metric on $\hat{\ov\BB}:=\pi^{-1}(\ov\BB)$. Then for any $0<s<1$, there exists $\la>0$ and a \Ka metric $\wt \omega$ on $\hat\CC^n$ so~that
\begin{equation*}
\wt{\omega}|_{\hat\BB_s}=\omega, \quad\wt{\omega}|_{\hat\CC^n\sminus \hat\BB}=\la\cdot\pi^*dd^c\rho.
\end{equation*}
\end{lemma}
\begin{proof}
There exists a pluriharmonic function $\psi$ on $\BB$, smooth and strictly plurisubharmonic on $\BB\sminus \{0\}$, so~that $\pi_*\omega=dd^c\psi$. By the maximum principle, the functions
\begin{equation*}
u,v:(0,1]\to \RR, \quad u(t)=\max_{\del\BB_t}\rho, \quad v(t)=\max_{\del\BB_t}\psi
\end{equation*}
are strictly increasing. It follows that for any $0<s<1$,
\begin{gather*}
r:=u(s)\equiv \rho|_{\del\BB_s}<R:=u(1)\equiv \rho|_{\del\BB}, \quad m:=\min_{\del\BB_s}\psi\leq v(s)<M:=v(1).
\end{gather*}
Thus, for any $\al>1$, we can take
\begin{equation*}
\la:=\al\,\frac{M-m}{R-r}>0, \quad c\in\left(M-\la R,m-\la r\right)\neq\emptyset,
\end{equation*}
so~that $\la\rho+c>\psi$ on $\del\BB$ and $\la\rho+c<\psi$ on $\del\BB_s$. Therefore, if $\wt{\psi}$ is the function on~$\CC^n$ defined as the regularized maximum between $\psi$ and $\la\rho+c$ on $\BB\sminus \BB_s$, equal to $\psi$ on $\BB_s$ and equal to $\la\rho+c$ on $\CC^n\sminus \BB$, then the metric $\pi^*(dd^c\wt{\psi})$ on $\hat\CC^n\sminus \pi^{-1}(0)$ glues to $\omega$ to define the desired metric $\wt{\omega}$ on $\hat\CC^n$.
\end{proof}

Now we are ready to characterize the existence of lcK metrics. For this, let us note that any proper modification $\pi:\hat\BB\to\BB$ induces naturally proper modifications $\pi:\hat\CC^n\to\CC^n$ and $\pi:\hat{\CC\PP}^n\to\CC\PP^n$, where $\hat{\CC\PP}^n$ is the compactification of $\hat\CC^n$ with a hyperplane at infinity.

\begin{theorem}\label{lcK}
Let $(\pi:\hat\BB\to\BB,\sigma:\overline{\BB}\to\hat\BB)$ be a Kato data and let $X=X(\pi,\sigma)$ be the corresponding Kato manifold.
The following are equivalent:
\begin{enumerate}
\item\label{lcK1} $X$ admits an lcK metric;
\item\label{lcK2} $\wt{X}$ admits a \Ka metric;
\item\label{lcK3} $\hat{\CC\PP}^n$ is a projective manifold;
\item\label{lcK4} $\hat{\mathbb{C}}^n$ admits a Kähler metric.

\end{enumerate}
\end{theorem}

\begin{proof}
The implications $\eqref{lcK1}\Rightarrow\eqref{lcK2}$ and $\eqref{lcK3}\Rightarrow\eqref{lcK4}$ are clear, while the implication $\eqref{lcK4}\Rightarrow \eqref{lcK1}$ is precisely Brunella's theorem \cite{bru} adapted to all complex dimensions \cite[Th.\,2.2]{iop}. Indeed, in \cite[Th.\,2.2]{iop} we showed that if $\pi$ is a composition of smooth blow-ups, then $X$ is lcK. The same proof works if one merely supposes that~$\hat\BB$, or equivalently $\hat\CC^n$, is Kähler. We are thus left with showing $\eqref{lcK2} \Rightarrow \eqref{lcK3}$.

If $\wt{X}$ is Kähler, then Proposition~\ref{embBall} implies that $\hat\BB\sminus \{\sigma(0)\}$ is also \Ka, so~by Miyaoka's extension theorem \cite[Prop.\,A]{miyaoka}, $\hat\BB$ admits a Kähler metric $\omega$. Now putting $\rho:=\log(1+\|z\|^2)$ in Lemma~\ref{extKA}, we find that $\hat\CC^n$ admits a \Ka metric $\wt{\omega}$ which glues to $\la\cdot\pi^*\omega_{FS}$ on $\hat{\CC\PP^n}\sminus \hat\BB$ for some $\la>0$, where $\omega_{FS}$ is the standard Fubini-Study metric on $\CC\PP^n$. Thus we have shown that $\hat{\CC\PP^n}$ is \Ka. Since on the other hand $\hat{\CC\PP^n}$ is a modification of $\CC\PP^n$, it is Moishezon, and therefore $\hat{\CC\PP^n}$ is projective by Moishezon's theorem \cite[Chap.\,I, Th.\,11]{moi}. This concludes the proof.
\end{proof}
